%%%%%%%%%%%%%%%%%%%%%%%%%%%%%%%%%%%%%%%%%%%%%%%%%%%%%%%%%%%%%%%%%%%%%%%
%
% Modelo de preenchimento de informações do cabeçalho
%
% Autor: Vitor G. P. Curtarelli
%
% Data: 10/03/2026
%
%%%%%%%%%%%%%%%%%%%%%%%%%%%%%%%%%%%%%%%%%%%%%%%%%%%%%%%%%%%%%%%%%%%%%%%

\titulo[Revisão bibliográfica do estado da arte]{Revisão bibliográfica do estado da arte de \textsl{Travel-Time Tomography}, inversão utilizando tempo de trânsito, e inversão do campo de velocidade do som na coluna d'água}  % Título do relatório % Argumento em [] (opcional) é o título curto, a ser usado hos headers (vide começo da pág. 2)
%\subtitulo{}  % Sub-título do relatório, se houver
\autor{Vitor G. P. Curtarelli}  % Autor(es) do relatório
% Vários autores podem ser adicionados em um único comando \autor, ou o comando pode ser chamado várias vezes
%O número de autores no cabeçalho (Autor/Autores) detecta e atualiza automaticamente quando houver 2+ autores
\revisor{Stephan Paul, Lucas C. Lobato} % Revisor(es) do relatório
% Uso idêntico ao comando \autor
\entregavel{3.1.1}{Revisão do SoA de técnicas e algoritmos para determinação de PVPS em coluna d'água}  % Número da atividade, como consta na Wiki do projeto
\resumo{Esse texto apresenta uma extensa revisão bibliográfica do estado da arte, tratando de diferentes tópicos relacionados à estimação do campo de velocidade do som na coluna d'água; com o propósito de elencar possíveis vias futuras de pesquisa a serem exploradas, que se alinhem com as limitações impostas pelo projeto. Também inclui uma introdução aos conceitos de estimação do campo de velocidade do som na água, bem como algumas limitações do projeto, e como essas se alinham com restrições de modelo ou matemáticas. }  % Um breve resumo do que o relatório trata